% !PRINTED 641
\weberpage{641}
\begin{center}
\S.\,176.\\[4pt]
\gesperrt{Der in $\Omega_m$ enthaltene reelle K\"orper $H_m$.}
\end{center}

Jede Zahl des K\"orpers $\Omega_m$ geht durch die Substitution $(r, r^{-1})$ in die conjugirt imagin\"are Zahl \"uber, die auch durch die Vertauschung $(i, -i)$ erhalten wird. Das Product zweier solcher conjugirt imagin\"arer Zahlen ist das \gesperrt{Quadrat des absoluten Werthes} einer jeden von ihnen.

Eine Zahl in $\Omega_m$ ist reell, wenn sie durch die Vertauschung $(r, r^{-1})$ unge\"andert bleibt, und nur unter dieser Voraussetzung. Jede solche Zahl l\"asst sich rational durch die \gesperrt{zweigliedrige Periode} $r + r^{-1}$ ausdr\"ucken, und geh\"ort also einem reellen K\"orper vom Grade $\frac12\varphi(m)$ an, der ein Theiler von $\Omega_m$ ist. Diesen wollen wir den in $\Omega_m$ \gesperrt{enthaltenen reellen K\"orper} nennen und mit $H_m$ bezeichnen (obwohl auch noch andere reelle K\"orper, n\"amlich alle Theiler von $H_m$, in $\Omega_m$ enthalten sind).

Die $\frac12\varphi(m) = \frac12\mu$ Zahlen
\weq{1}{1,\ r+r^{-1},\ r^2+r^{-2},\ \dots,\ r^{\frac12\mu-1}+r^{-\frac12\mu+1}}
bilden eine Basis der ganzen Zahlen des K\"orpers $H_m$. Denn nach \S.\,170, (7) ist
\weq{2}{\omega = x_0+x_1r+x_2r^2+\cdots+x_{\frac12\mu-1}r^{\frac12\mu-1}+x_{\frac12\mu}r^{\frac12\mu} +x_{-1}r^{-1}+x_{-2}r^{-2}+\cdots+x_{-\frac12\mu+1}r^{-\frac12\mu+1}}
dann und nur dann eine ganze Zahl des K\"orpers $\Omega_m$, wenn die Co\"efficienten $x_0, x_1, \dots$ ganze rationale Zahlen sind.

Die Bedingung daf\"ur, dass diese Zahl dem K\"orper $H_m$ angeh\"ort, erh\"alt man, wenn man die Substitution $(r, r^{-1})$ ausf\"uhrt und die Differenz $=0$ setzt. Dividirt man noch durch $r-r^{-1}$, so ergiebt sich so:
\[
(x_1-x_{-1})+(x_2-x_{-2})(r+r^{-1})+\cdots
+(x_{\frac12\mu-1}-x_{-\frac12\mu+1})\frac{r^{\frac12\mu-1}-r^{-\frac12\mu+1}}{r-r^{-1}}
+x_{\frac12\mu}\frac{r^{\frac12\mu}-r^{-\frac12\mu}}{r-r^{-1}}=0.
\]
Die Divisionen lassen sich hier ausf\"uhren, und wenn man dann mit $r^{\frac12\mu-1}$ multiplicirt, so ergiebt sich f\"ur $r$ eine rationale Gleichung von niedrigerem als $\mu^{\text{ten}}$ Grade, die nur dann be-
%CONT>>
